\section{Naturalidad conjunta de las publicaciones del estado genealógico}
\label{sec:cpe-naturalidad-tres-publicaciones}

Sea \((X_m,p_{m+1,m})_{m\ge0}\) el sistema inverso de estados
APP--TRIT--TPK calibrados. En cada nivel se conservan fibra, relación,
número, orientación, acarreo, memoria, frontera, ruta, multisección,
supervivencia, terminal y lector. Sean
\[
 \mathfrak R_m:X_m\to A_m,
 \qquad
 \mathfrak E_{m,\chi}:X_m\to E_m,
 \qquad
 \mathfrak D_m:X_m\to D_m
\]
las publicaciones arquimediana, excepcional calibrada y dimensional. Sus
truncamientos se denotan por \(a_{m+1,m}\), \(e_{m+1,m}\) y
\(d_{m+1,m}\).

\begin{theorem}[Cono natural de las \(3\) publicaciones]
\label{thm:cpe-cono-tres-publicaciones}
Si los lectores de nivel finito satisfacen
\begin{equation}
\begin{aligned}
 a_{m+1,m}\mathfrak R_{m+1}&=\mathfrak R_mp_{m+1,m},\\
 e_{m+1,m}\mathfrak E_{m+1,\chi}&=\mathfrak E_{m,\chi}p_{m+1,m},\\
 d_{m+1,m}\mathfrak D_{m+1}&=\mathfrak D_mp_{m+1,m},
\end{aligned}
\label{eq:cpe-naturalidad-tres-lectores}
\end{equation}
entonces existen y son únicos los lectores de límite
\[
 \mathfrak R_\infty:X_\infty\to A_\infty,
 \quad
 \mathfrak E_{\infty,\chi}:X_\infty\to E_\infty,
 \quad
 \mathfrak D_\infty:X_\infty\to D_\infty
\]
que reproducen todos los niveles. El mapa conjunto
\begin{equation}
 \mathfrak J_\chi
 :=(\mathfrak R_\infty,\mathfrak E_{\infty,\chi},\mathfrak D_\infty)
 :X_\infty\longrightarrow A_\infty\times E_\infty\times D_\infty
 \label{eq:cpe-mapa-conjunto-tres-publicaciones}
\end{equation}
es una publicación del mismo estado y no una concatenación de tres
generadores.
\end{theorem}

\begin{proof}
Para \(x=(x_m)_m\in X_\infty\), las tres sucesiones
\((\mathfrak R_mx_m)_m\), \((\mathfrak E_{m,\chi}x_m)_m\) y
\((\mathfrak D_mx_m)_m\) son compatibles por
\eqref{eq:cpe-naturalidad-tres-lectores}; definen, por la propiedad universal
del límite inverso, los tres lectores. La misma propiedad prueba su unicidad.
La fórmula \eqref{eq:cpe-mapa-conjunto-tres-publicaciones} evalúa una sola
historia \(x\) y por tanto conserva el antecedente causal común.
\end{proof}

La igualdad legítima no afirma que una salida con pérdida reconstruya por sí
sola todas las demás. Afirma que los tres cuadrados de truncación conmutan y
que sus salidas son coordenadas naturales de un mismo cono. En particular,
la realización arquimediana puede olvidar genealogía sin que ese olvido se
atribuya a la publicación excepcional o a la dimensional.

\section{Separación conjunta de valor e incidencia excepcional}
\label{sec:cpe-separacion-conjunta}

Fijado el calibre \(\chi_{\rm exc}\), el lector excepcional actúa sobre cada
generador por
\begin{equation}
 (w_j,f_j)\longmapsto
 \bigl(f_j,\Gamma_{f_jA_Wf_j^{-1}}(w_j)\bigr).
 \label{eq:cpe-lector-excepcional-generador}
\end{equation}
El marco \(f_j\) forma parte de la salida y
\(\Gamma_{f_jA_Wf_j^{-1}}\) es un automorfismo de la fibra calibrada.

\begin{theorem}[Separación calibrada de historias]
\label{thm:cpe-separacion-calibrada}
El lector excepcional
\(\mathfrak E_{\infty,\chi_{\rm exc}}\) es inyectivo sobre el cociente por
los cambios de coordenadas incluidos en \(\chi_{\rm exc}\). En consecuencia,
\[
 (\mathfrak R_\infty,\mathfrak E_{\infty,\chi_{\rm exc}})
 :X_\infty^{\rm cal}\longrightarrow
 A_\infty\times E_\infty
\]
es un encaje topológico cuando \(X_\infty^{\rm cal}\) es compacto y el
codominio es de Hausdorff.
\end{theorem}

\begin{proof}
De la salida \((f_j,y_j)\), donde
\(y_j=\Gamma_{f_jA_Wf_j^{-1}}(w_j)\), se recupera
\[
 w_j=\Gamma_{f_jA_Wf_j^{-1}}^{-1}(y_j).
\]
Esto reconstruye cada generador y, por compatibilidad de truncación, toda la
historia. Dos representantes que difieren sólo por un cambio incluido en el
calibre definen la misma clase; fuera de esa equivalencia, la reconstrucción
es única. La inyectividad conjunta es inmediata. Una aplicación continua e
inyectiva de un compacto a un Hausdorff es un homeomorfismo sobre su imagen.
\end{proof}

\section{Antecedente no conmutativo común de la doble proyección}
\label{sec:cpe-antecedente-ucp}

Sea
\(\mathcal A_{\rm TPK}=C(X_F^{\leftrightarrow})\rtimes_\sigma\mathbb Z\)
y sea \(u\) el unitario que implementa \(\sigma\). Los lectores continuos
\(\mathfrak R\) y \(\mathfrak E_\chi\) inducen los morfismos unitales
\begin{equation}
 \Phi_{\rm arq}(f)=f\circ\mathfrak R,
 \qquad
 \Phi_{\rm exc}(g)=g\circ\mathfrak E_\chi,
 \label{eq:cpe-dos-mapas-ucp}
\end{equation}
desde las álgebras conmutativas de sus codominios hacia
\(C(X_F^{\leftrightarrow})\subset\mathcal A_{\rm TPK}\).

\begin{theorem}[Sistema de operadores común y equivarianza TPK]
\label{thm:cpe-sistema-operadores-comun}
El subespacio autoadjunto unital
\[
 \mathcal S_\infty
 =\operatorname{span}\bigl\{
 I,u,u^*,\Phi_{\rm arq}(C(A_\infty)),
 \Phi_{\rm exc}(C(E_\infty))\bigr\}
 \subset\mathcal A_{\rm TPK}
\]
es un sistema de operadores no conmutativo. Los dos mapas de
\eqref{eq:cpe-dos-mapas-ucp} son UCP\@. Si los lectores entrelazan las
dinámicas de salida \(\alpha_{\rm arq}\) y \(\alpha_{\rm exc}\), entonces
\begin{equation}
 \operatorname{Ad}_u\Phi_{\rm arq}
 =\Phi_{\rm arq}\alpha_{\rm arq},
 \qquad
 \operatorname{Ad}_u\Phi_{\rm exc}
 =\Phi_{\rm exc}\alpha_{\rm exc}.
 \label{eq:cpe-equivarianza-doble-ucp}
\end{equation}
\end{theorem}

\begin{proof}
La composición con un mapa continuo es un homomorfismo unital de
\(C^*\)-álgebras y, por tanto, es completamente positiva. El unitario
\(u\) hace no conmutativa la envolvente cuando la dinámica es no trivial.
La identidad \eqref{eq:cpe-equivarianza-doble-ucp} se verifica evaluando
ambos miembros en una historia y usando la equivarianza de cada lector.
\end{proof}

Esta construcción no fuerza dos mediciones qutrit mutuamente insesgadas a
poseer una PVM conjunta. El antecedente común es el álgebra cruzada de
historias; las publicaciones son dos canales UCP distintos desde sus
codominios.

